\documentclass[a4paper,landscape,11pt]{article}

\usepackage[utf8]{inputenc}
\usepackage[T2A]{fontenc}
\usepackage[russian]{babel}
\usepackage{geometry}
\usepackage{tikz}
\usetikzlibrary{arrows.meta, positioning}

\pagestyle{empty}
\geometry{margin=12mm}

\tikzset{
	uml/.style={
		draw,
		rounded corners=1pt,
		align=left,
		font=\scriptsize,
		text width=43mm,
		inner sep=3pt,
		minimum height=12mm
	},
	relation/.style={
		-{Stealth[length=2mm]},
		line width=.45pt
	},
	dependency/.style={
		-{Stealth[length=2mm]},
		dashed,
		line width=.45pt
	}
}

\begin{document}
	
	\begin{center}
		\Large\textbf{UML-диаграмма классов клиент-серверного приложения}
	\end{center}
	
	\vspace{4mm}
	\centering
	
	\begin{tikzpicture}[x=1cm, y=1cm]
		
		% ===================== Клиентская часть =====================
		
		\node[uml] (ClientMain) at (-10,0)
		{
			\textbf{client::Main}\\
			\hrulefill\\
			+ main(String[]): void
		};
		
		\node[uml] (ClientConsole) at (-10,3)
		{
			\textbf{client::Console}\\
			\hrulefill\\
			+ readln(): String\\
			+ println(Object): void\\
			+ printError(Object): void
		};
		
		\node[uml] (ClientCommandManager) at (-4,3)
		{
			\textbf{client::CommandManager}\\
			\hrulefill\\
			- client: DatagramClient\\
			- commands: Map<String, Command>\\
			- activeScripts: Set<Path>\\
			- activeAsks: Deque<Ask>
		};
		
		\node[uml] (DatagramClient) at (-4,0)
		{
			\textbf{client::DatagramClient}\\
			\hrulefill\\
			- channel: DatagramChannel\\
			- server: SocketAddress\\
			+ request(Request): Response\\
			+ close(): void
		};
		
		\node[uml] (ClientCommand) at (2,3)
		{
			\textbf{client::Command}\\
			\hrulefill\\
			+ type(): CommandType\\
			+ apply(Request): Response\\
			\hrulefill\\
			Help, Info, Show, Add,\\
			Update, Save, Exit, \ldots
		};
		
		% ===================== Общие классы =====================
		
		\node[uml] (DatagramTransfer) at (-4,-4)
		{
			\textbf{common::DatagramTransfer}\\
			\hrulefill\\
			+ send(channel, data, addr): void\\
			+ receive(channel, timeout): ReceivedData\\
			+ receiveAvailable(...): ReceivedData
		};
		
		\node[uml] (request) at (-10,-4)
		{
			\textbf{common::Protocol.Request}\\
			\hrulefill\\
			- command: CommandType\\
			- payload: Payload
		};
		
		\node[uml] (response) at (2,-4)
		{
			\textbf{common::Protocol.Response}\\
			\hrulefill\\
			- ok: boolean\\
			- message: String
		};
		
		\node[uml] (Model) at (8,-4)
		{
			\textbf{common::Model.SpaceMarine}\\
			\hrulefill\\
			- id, name, health\\
			- coordinates: Coordinates\\
			- creationDate: ZonedDateTime\\
			- chapter: Chapter\\
			+ compareTo(SpaceMarine)
		};
		
		% ===================== Серверная часть =====================
		
		\node[uml] (ServerMain) at (-10,-9)
		{
			\textbf{server::Main}\\
			\hrulefill\\
			+ main(String[]): void
		};
		
		\node[uml] (ServerConsole) at (-10,-12)
		{
			\textbf{server::Console}\\
			\hrulefill\\
			+ readln(): String\\
			+ println(Object): void\\
			+ printError(Object): void
		};
		
		\node[uml] (datagram) at (-4,-9)
		{
			\textbf{server::DatagramServer}\\
			\hrulefill\\
			- channel: DatagramChannel\\
			+ run(): void\\
			+ stop(): void
		};
		
		\node[uml] (handler) at (-4,-12)
		{
			\textbf{managers::ServerCommandManager}\\
			\hrulefill\\
			+ handle(Request): Response\\
			+ runConsole(): void
		};
		
		\node[uml] (command) at (2,-9)
		{
			\textbf{commands::Command}\\
			\hrulefill\\
			+ type(): CommandType\\
			+ apply(Request): Response\\
			\hrulefill\\
			Add, Update, RemoveById,\\
			Save, Exit, \ldots
		};
		
		\node[uml] (service) at (8,-9)
		{
			\textbf{managers::CollectionManager}\\
			\hrulefill\\
			- collection: LinkedList<SpaceMarine>\\
			- currentId: int\\
			+ add(), update(), clear()\\
			+ snapshot(), descending()
		};
		
		\node[uml] (storage) at (8,-12)
		{
			\textbf{managers::DumpManager}\\
			\hrulefill\\
			+ readCollection(): LinkedList\\
			+ writeCollection(): void
		};
		
		\node[uml] (manager) at (2,-12)
		{
			\textbf{managers::CommandManager}\\
			\hrulefill\\
			- commands: Map\\
			- history: Deque\\
			+ register(), getCommand()
		};
		
		% ===================== Связи: клиент =====================
		
		\draw[dependency]
		(ClientMain) -- node[above, font=\scriptsize]{создаёт} (DatagramClient);
		
		\draw[dependency]
		(ClientMain) -- node[above, font=\scriptsize]{создаёт} (ClientCommandManager);
		
		\draw[dependency]
		(ClientMain) -- node[left, font=\scriptsize]{создаёт} (ClientConsole);
		
		\draw[dependency]
		(ClientCommandManager) -- node[above, font=\scriptsize]{использует} (ClientConsole);
		
		\draw[dependency]
		(ClientCommandManager) -- node[left, font=\scriptsize]{использует} (DatagramClient);
		
		\draw[dependency]
		(ClientCommandManager) -- node[above, font=\scriptsize]{использует} (ClientCommand);
	
		
		% ===================== Связи: сервер =====================
		
		\draw[dependency]
		(ServerMain) -- node[left, font=\scriptsize]{создаёт} (handler);
		
		\draw[dependency]
		(ServerMain) -- node[above, font=\scriptsize]{запускает} (datagram);
		
		\draw[dependency]
		(ServerMain) -- node[left, font=\scriptsize]{создает} (ServerConsole);
		
		\draw[dependency]
		(handler) -- node[above, font=\scriptsize]{использует} (ServerConsole);
		
		\draw[dependency]
		(datagram) -- node[right, font=\scriptsize]{вызывает} (handler);
		
		\draw[dependency]
		(handler) -- node[above, font=\scriptsize]{выбирает} (command);
		
		\draw[dependency]
		(handler) -- node[above, font=\scriptsize]{использует} (manager);
		
		\draw[dependency]
		(manager) -- node[below, font=\scriptsize]{хранит} (command);
		
		\draw[relation]
		(command) -- node[above, font=\scriptsize]{использует} (service);
		
		\draw[dependency]
		(storage) -- node[right, font=\scriptsize]{загружает/сохраняет} (service);
		
		% ===================== Связи: UDP-обмен =====================
		
		\draw[dependency]
		(datagram) -- node[right, font=\scriptsize]{использует} (DatagramTransfer);
		
		\draw[dependency]
		(DatagramClient) -- node[right, font=\scriptsize]{использует} (DatagramTransfer);
		
		\draw[dependency]
		(datagram) -- node[right, font=\scriptsize]{десериализует} (request);
		
		\draw[dependency]
		(datagram) -- node[right, font=\scriptsize]{сериализует} (response);
		
		\draw[dependency]
		(DatagramClient) -- node[above, font=\scriptsize]{сериализует} (request);
		
		\draw[dependency]
		(DatagramClient) -- node[right, font=\scriptsize]{десериализует} (response);

		
	\end{tikzpicture}
	
\end{document}